% 第八章 結論與未來工作
\chapter{結論與未來工作}\label{ch:conclusion}

\section{結論}\label{sec:conclusion}

本專題在兩家廠商的現役 GPU（NVIDIA GH200 叢集與兩個 8 張 AMD MI350X 的節點）上，以純軟體將用不到的尾數位元歸零，回答了第 \ref{sec:questions} 節的四個問題：

\begin{enumerate}
  \item \textbf{機制：}在兩個平台的矩陣引擎、FP64 向量單元與記憶體搬移上，以全 1 或固定樣式填入被捨去的位元，都保留了大部分的效益（41–97\%），所以省電主要來自低位元不再翻轉，數值為 0 再多省一些。在數值與 1 的個數都固定時，只改變運算元的排列順序就讓每次運算的能量相差 8–9 個百分點，相鄰運算元之間的翻轉數解釋 83\% 的變異，所以主因是相鄰運算元之間變化的位元。MI350X 上，截斷在鎖定時脈下還提高了有效時脈，可能來自電路層級的時脈延展。
  \item \textbf{效益：}在功耗受限且截斷實作於程式內時，以矩陣乘法為主、具正確性檢查的完整程式在兩個平台上都獲得正面收益。GH200 上，HPL（調參後）保留 36 位有效位數在 500–800~W 快 3.3–4.7\%，HPL-MxP 保留 3 位在 500~W 快 5.8\%，Quantum ESPRESSO 依收斂門檻調整位元數，在總能量印出的每一位都不變的情況下快 2.1–3.6\%，DF-MP2 在 500–700~W 快 4.7–7.7\%，高斯過程快 7.5–7.9\%；在預設的 1000~W 主要表現為省電（HPL 每次求解少耗電 7.5–8.6\%）。MI350X 的預設上限就壓住時脈，所以在出廠設定下 HPL 快 $4.8\pm0.3$\%、HPL-MxP $5.0\pm0.4$\%、高斯過程 $7.4\pm0.6$\%、DF-MP2 $10.3\pm0.6$\%，在 500~W 快 7.0–19.3\%（兩節點合併的分析單位）。HPL-MxP 的結果只適用於對角優勢的測試矩陣。代價是必須在每個平台上以多組問題大小驗證保留位數（HPL 28 位在 GH200 大多通過、在 MI350X 全部失敗），並使用最近捨入。
  \item \textbf{預測：}若截斷只改變切換電容，截斷版本在上限 $P$ 下就等於原生版本在 $P_0+(P-P_0)/R$ 下。以實測的閒置功耗作為 $P_0$，這個模型對兩個平台各 42 個矩陣乘法的點平均誤差為 1.5 與 2.1 個百分點，對兩個沒有參與建模的應用為 1.6 個百分點（最大 3.6）。因此只要知道平台的閒置功耗與應用的原生曲線，一次鎖頻量測就能估計截斷版本在每個上限下的效能，省去大部分的上限掃描；但在兩種方法只差一兩個百分點的交叉點附近，這樣的誤差還不足以可靠地選出較快的方法。
  \item \textbf{適用範圍：}尾數歸零與 INT8 模擬 FP64 互補。MI350X 上，模擬在內積長、上限寬鬆時大幅領先（DF-MP2 快 67\%；HPL 在它本身的殘差門檻下，1000~W 快 34\%、500~W 快 12\%，都勝過截斷）；把共同的殘差門檻收緊到 1（這是敏感度分析，不是 HPL 的要求），500~W 下已測設定的排序反轉（截斷 +7.3\%，模擬 $-4.0$\%）。模擬在內積短時比原生還慢，需要的工作空間迫使 HPL 縮小問題，在列內指數跨度寬時會安靜地損失精度（GH200 上的高斯基底 DFT 錯 52.5\mEh）；我們所用的 GEMMul8 HIP 版本另有一個尺寸相關的實作缺陷。把所有 GEMM 換成 FP32/3xTF32 在長內積上無法取代截斷；在已使用 FP8 的 AI 上，FP8 已取得幾乎全部的好處。本報告比較的兩種方法中，只有截斷適用於 BF16 與 FP8。 只用矩陣乘法量測擬合的逐次呼叫選擇規則，在完全沒有參與擬合的 DF-MP2 上 99.3\% 的呼叫選對，只比逐次最佳慢 0.02\%。
\end{enumerate}

另外，在功耗受限的量測中，溫度是必須控制的干擾：MI350X 的機箱風扇依前一次執行調整，使下一次執行快多達 4.5\%，交錯順序無法消除；我們以溫度迴歸校正，並以穩態與相同的熱歷史兩種受控實驗確認了校正後的數值。

總結而言，在功耗受限的 FP64 科學計算中，尾數歸零是一個每個元素只多一次位元運算、不需要新硬體、精度可以 1 位為單位調整、而且能事先預測收益的調整手段，最適合具殘差修正或自洽迭代的求解器。隨著吞吐量成長得比功耗預算快，在上限下運作會成為常態，翻轉用不到的位元所消耗的能量也會成為效能的一階參數。

\section{未來工作}\label{sec:future}

\begin{description}
  \item[更多硬體：] 在 FP64 相對弱的平台（例如 Blackwell 的消費級或資料中心 GPU）與其他世代上重做分界掃描與 HPL，檢驗模型與分界能否移植；在多節點 MI350X 上重做 HPL 與 HPL-MxP。
  \item[正式規則下的 HPL-MxP：] 改用規則允許的矩陣產生器，重做保留位數的掃描。
  \item[函式庫整合：] 把「依形狀、上限與列內跨度選擇截斷或模擬」的規則實作於矩陣乘法函式庫，使截斷不需要額外複製；功耗模型可以提供上限這一側的判斷。
  \item[更多自適應求解器：] 把「依容許誤差決定位元數」推廣到其他迭代法，例如 Krylov 求解器的內層與混合精度迭代修正；也可以利用資料的位元資訊量事先估計需要保留的位數~\cite{Klower2021}，以減少試驗。
  \item[AI 的第二項研究：] 在功耗上限下，降低精度能否縮短達到目標品質的時間與能耗，需要另外設計實驗（附錄 \ref{app:ai} 只是初步的負面結果）。
  \item[量測：] 以外接電表交叉驗證絕對能量；在可以固定風扇轉速的環境下重做溫度控制實驗。
\end{description}
